\documentclass[12pt,a4paper]{article}
\usepackage[utf8]{inputenc}
\usepackage[czech]{babel}
\usepackage[T1]{fontenc}
\usepackage{hyperref}
\usepackage{graphicx}
\usepackage{listings}
\usepackage{xcolor}
\usepackage{geometry}
\usepackage{tabularx}

\geometry{
    a4paper,
    total={170mm,257mm},
    left=20mm,
    top=20mm,
}

\hypersetup{
    colorlinks=true,
    linkcolor=black,
    filecolor=magenta,      
    urlcolor=blue,
}

\renewcommand{\maketitle}{%
 \begin{titlepage}
  \centering
  \vspace*{3cm}
  {\Huge\bfseries KIV/WEB}\par
  \vspace{0.5cm}
  {\Large Semestrální práce}\par
  \vspace{1.5cm}
  {\Large \textbf{Konferenční systém}}\par
  \vspace{0.5cm}
  {Webová aplikace pro správu konference}\par
  \vfill
  \begin{flushleft}
   \textbf{Autor:} Marek Jirušek\\
   \textbf{Email:} marjir@students.zcu.cz\\ % DOPLŇTE EMAIL
   \textbf{Datum:} \today
  \end{flushleft}
 \end{titlepage}%
}

\begin{document}

\maketitle
\newpage

\tableofcontents
\newpage

\section{Zadání a popis aplikace}

Cílem semestrální práce bylo navrhnout a implementovat webovou aplikaci pro  správu odborné konference. Systém slouží k registraci účastníků, odevzdávání odborných příspěvků (článků) ve formátu PDF a k zajištění celého procesu recenzního řízení.



\subsection{Funkcionalita podle rolí}
Aplikace rozlišuje několik uživatelských rolí, z nichž každá má specifická oprávnění:

\begin{itemize}
    \item \textbf{Neregistrovaný návštěvník:} Může prohlížet úvodní stránku s nejnovějšími články, procházet archiv schválených příspěvků a provést registraci či přihlášení.
    \item \textbf{Čtenář:} Základní role po registraci. Má stejná práva jako návštěvník, ale je evidován v systému.
    \item \textbf{Autor:} Může vytvářet nové články, nahrávat k nim PDF soubory, upravovat je (pokud jsou ve stavu konceptu nebo vrácené k přepracování) a odesílat je do recenzního řízení.
    \item \textbf{Recenzent:} Vidí seznam článků, které mu byly administrátorem přiděleny k recenzi. Může článek přijmout či zamítnout a musí k němu vypracovat slovní hodnocení.
    \item \textbf{Administrátor:} Má na starosti správu uživatelů (změna rolí, mazání účtů) a správu recenzního procesu (přiřazování recenzentů k článkům). Vidí stav recenzí a na jejich základě článek finálně schvaluje (publikuje) nebo zamítá.
    \item \textbf{Super-Admin:} Speciální role s nejvyšším oprávněním, která nemůže být běžným administrátorem smazána ani upravena.
\end{itemize}

\section{Použité technologie}

Aplikace je postavena na architektuře \textbf{MVC} (Model-View-Controller) 
\subsection{Backend}
\begin{itemize}
    \item \textbf{PHP 8.x:} Hlavní serverový jazyk zajišťující logiku aplikace, routing a komunikaci s databází.
    \item \textbf{MySQL / MariaDB:} Relační databáze pro ukládání dat (uživatelé, články, recenze).
    \item \textbf{PDO (PHP Data Objects):} Použito pro přístup k databázi. Zajišťuje bezpečnost pomocí \textit{Prepared Statements} (ochrana proti SQL Injection).
\end{itemize}

\subsection{Frontend}
\begin{itemize}
    \item \textbf{HTML5 / CSS3:} Sémantické značkování obsahu. Vlastní styly jsou umístěny v adresáři \texttt{CSSs/} a doplňují framework.
    \item \textbf{Bootstrap 5:} CSS framework využitý pro zajištění responzivního designu (Grid systém) a pro UI komponenty (navigační lišta, modální okna, alerty, karty).
    \item \textbf{JavaScript (ES6):}
    \begin{itemize}
        \item Využit pro asynchronní komunikaci (\textbf{AJAX}) v sekci pro správu recenzentů (přiřazování bez obnovení stránky).
        \item Obsluha UI prvků (přepínání mezi zobrazením PDF a textového obsahu).
    \end{itemize}
    \item \textbf{CKEditor 5:} Externí knihovna (WYSIWYG editor) použitá pro formátování textového obsahu článků.
\end{itemize}

\section{Adresářová struktura}

Struktura projektu je rozdělena logicky podle odpovědností MVC architektury:

\begin{itemize}
    \item \textbf{config/} -- Konfigurační soubory (připojení k DB, konstanty).
    \item \textbf{Controlers/} -- Třídy řadičů (Controllers), které zpracovávají požadavky.
    \item \textbf{CSSs/} -- Kaskádové styly (vlastní CSS soubory).
    \item \textbf{Models/} -- Třídy pro práci s daty a databází.
    \item \textbf{uploads/} -- Úložiště pro nahrané PDF soubory článků.
    \item \textbf{Views/} -- HTML šablony (pohledy) pro zobrazení obsahu uživateli.
    \item \textbf{index.php} -- Vstupní bod aplikace (Front Controller).
    
\end{itemize}

\section{Architektura aplikace}

Aplikacevyužívá objektově orientovaný přístup (OOP). Vstupním bodem je \texttt{index.php}, který inicializuje \texttt{Routing}.

\subsection{Popis klíčových tříd}

\subsubsection{Routing a Database modely}
\begin{itemize}
    \item \textbf{\texttt{Routing}:} Zpracovává URL parametry (např. \texttt{?path=article-maker}), na jejichž základě dynamicky načítá příslušný Controller a předává mu řízení.
    \item \textbf{\texttt{Database}:} Třída implementující návrhový vzor \textit{Singleton}. Zajišťuje jediné připojení k databázi pro celou dobu běhu skriptu.
\end{itemize}

\subsubsection{Modely}
Modely zapouzdřují logiku pro práci s daty a komunikaci s databází.
\begin{itemize}
    \item \textbf{\texttt{UserModel}:} Zajišťuje autentizaci (ověření hesla), registraci a správu uživatelů.
    \item \textbf{\texttt{ArticleModel}:} Implementuje operace pro články. Vrací data jako instance třídy \texttt{Article}.
    \item \textbf{\texttt{ReviewManagementModel}:} Obsahuje logiku pro administraci recenzí – přiřazování recenzentů a změny stavů článků.
    \item \textbf{\texttt{RecenzentModel}:} Poskytuje data specifická pro roli recenzenta (seznam přidělených článků, ukládání hodnocení).
\end{itemize}

\subsubsection{Controllery (Controllers)}
Controllery propojují Modely a Views. Ověřují oprávnění uživatele.
\begin{itemize}
    \item \textbf{\texttt{ArticleMakerController}:} Řídí proces vytváření a editace článku autorem, včetně uploadu souborů a validace.
    \item \textbf{\texttt{ManageUsersController}:} Zajišťuje výpis uživatelů a změny jejich rolí.
    \item \textbf{\texttt{ManageRecenzentsController}:} Zpracovává požadavky na přiřazení recenzentů.
    \item \textbf{\texttt{ArticleController}:} Zobrazuje detail článku a zajišťuje bezpečné stažení PDF (kontrola oprávnění).
\end{itemize}

\section{Seznam defaultních uživatelů}

V databázi jsou předpřipraveni uživatelé pro otestování všech rolí aplikace.

\begin{center}
\renewcommand{\arraystretch}{1.5}
\begin{tabular}{|l|l|l|l|}
\hline
\textbf{Login} & \textbf{Heslo} & \textbf{Role} & \textbf{Poznámka} \\ \hline
SUPERADMIN & SUPERADMIN & Super-Admin & Stejná práva jako admin, může přidělit admina \\ \hline
LOL & LOLLOL & Admin & Správa uživatelů a recenzí \\ \hline
NOOB & NOOBNOOB & Recenzent & Hodnotí články \\ \hline
KOŠ & KOŠKOŠ & Recenzent & Hodnotí články \\ \hline
TEST & TESTTEST & Autor & Vytváří články \\ \hline
PEPA & PEPAPEPA & Čtenář & Pouze prohlížení \\ \hline
IDK & IDKIDK & Recenzent & Hodnotí články \\ \hline
\end{tabular}
\end{center}

\end{document}